\section{Exterior channels and the density near the median plane}
\label{sec:channels}

The leading removal amplitude has a simple zero on the median plane,
but a complementary ionic channel supplies a nonzero density tail
there with the same exponential rate.  This is the nodal-plane
mechanism discussed in \cite[Section~V.A]{GGB}.  We prove the estimates,
including their uniformity as the interaction tends to zero, without
assuming an asymptotic expansion of finitely many Dyson orbitals.

In the perturbative part of this section \(D\) is fixed and sufficiently
large, and \(0\leq\lambda\leq\lambda_D\), where \(\lambda_D>0\) may be
decreased finitely many times.  Exterior radii and constants denoted
by \(C_D\) may depend on \(D\), but not on \(\lambda\) in this interval.
For \(x=(x_1,x_2,z)\), write \(r=|x|\),
\(p=(x_1^2+x_2^2)^{1/2}\), and \(\omega=x/r\) when \(r>0\).
Differentiation in \(p\) keeps \(z\) fixed.  Put
\begin{equation}
 k=\sqrt{2(E_2(\lambda)-E_3(\lambda))},\qquad
 \beta=\frac{2-2\lambda}{k}-1,\qquad A(r)=e^{-kr}r^\beta.
 \label{eq:channel-radial-factor}
\end{equation}
We can arrange that
\begin{equation}
 k>0.99,\qquad \gamma-k\geq\tfrac12(\gamma-k_0)>0.
 \label{eq:channel-rate-separation}
\end{equation}

\begin{proposition}[Density estimates]
\label{prop:density-properties}
For all sufficiently large \(D\) and sufficiently small \(\lambda\geq0\),
the density \(n=n_\lambda\) is positive, even in \(z\), and axisymmetric.
It is smooth away from the nuclei and converges locally uniformly to
\(2b^2+e^2\) as \(\lambda\downarrow0\).

Given \(\eta>0\), first \(D\) can be chosen sufficiently large and then
\(\lambda_D>0\) sufficiently small so that, beyond a radius independent
of \(0\leq\lambda\leq\lambda_D\),
\begin{equation}
 |\partial_r\log n+2k|\leq\eta,\qquad
 \left|\partial_p\log n+2k\frac p r\right|\leq\eta
 \quad(0\leq z\leq p/2).
 \label{eq:density-drift}
\end{equation}
With \(L=10\) and \(m(p,z)=e^{2kp}n(p,z)\), the same choices give
\begin{equation}
 \frac{|m_p|+|m_{pp}|}{m}\leq\eta,\qquad
 \frac{n(p,z_1)}{n(p,z_2)}\leq20
 \quad(0\leq z_1\leq z_2\leq L)
 \label{eq:strip-data}
\end{equation}
for all sufficiently large \(p\), uniformly in that coupling interval.

For each fixed such \(\lambda>0\), there are \(c_\lambda>0\) and
\(P_\lambda<\infty\) such that
\begin{equation}
 \sqrt{n(p,0)}\geq c_\lambda A(p)p^{-4}\qquad(p\geq P_\lambda).
 \label{eq:plane-tail}
\end{equation}
There is a number \(a>0\), depending on \(D,\lambda\), with the following
translation properties.  If \(x_j=p_j\omega_j\) lies in the median plane,
\(p_j\to\infty\), and the horizontal unit vectors \(\omega_j\to\omega_*\),
then
\begin{equation}
 \frac{n(x_j+x)}{(A(p_j)/p_j)^2}
 \longrightarrow a^2z^2e^{-2k\omega_*\cdot x}
 \quad\hbox{locally uniformly for }x=(x_1,x_2,z)\in\R^3.
 \label{eq:plane-translate}
\end{equation}
If \(|x_j|\to\infty\), \(|(x_j)_3|\to\infty\), and
\(x_j/|x_j|\to\omega_*\), then
\begin{equation}
 \frac{n(x_j+x)}{n(x_j)}
 \longrightarrow e^{-2k\omega_*\cdot x}
 \quad\hbox{locally uniformly in }\R^3.
 \label{eq:bulk-translate}
\end{equation}
\end{proposition}

The quantitative uniformity needed for the final parameter prescription
is given in \Cref{eq:quantitative-density-drift} below.  In particular,
the constants remaining after the exterior limit do not involve the
inverse of the small splitting \(d_*\).

\subsection{The fiber equation and an exterior cutoff}

Let \(\cK\) be the Hilbert space of the antisymmetric ionic pair,
including its spin coordinates and the distinguished electron's spin,
restricted to total spin projection \(1/2\).  Regard
\[
 F_{\mathrm{true}}(x)=\sqrt3\,\Psi_\lambda(x,\cdot)
\]
as a \(\cK\)-valued function, with \(\|F_{\mathrm{true}}(x)\|^2=n(x)\).
The ionic operator acts trivially on the distinguished spin.  Set
\[
 B=H_{\mathrm{ion}}-E_2(\lambda),\qquad
 P=|\phi\rangle\langle\phi|,\qquad Q=1-P,\qquad B_Q=B|_{Q\cK},
\]
where \(\phi\) is the normalized ionic singlet tensored with distinguished
spin up.  Only that distinguished spin is compatible with the singlet
and total projection \(1/2\); hence \(P\) has rank one and
\begin{equation}
 B\geq0,\qquad B_Q\geq c_D>0
 \label{eq:ionic-gap}
\end{equation}
uniformly on the small coupling interval.

For normalized spin orbitals \(v,w\), let \(\Phi_{v,w}\) denote their
normalized two-electron determinant.  Expansion of the unperturbed
determinant in the distinguished coordinate gives
\begin{equation}
 \begin{split}
 F_0(x)&=e(x)\phi_0+b(x)G,\\
 \phi_0&=\uparrow\otimes\Phi_{b\uparrow,b\downarrow},\\
 G&=\uparrow\otimes\Phi_{b\downarrow,e\uparrow}
       -\downarrow\otimes\Phi_{b\uparrow,e\uparrow}.
 \end{split}
 \label{eq:unperturbed-fibers}
\end{equation}
Consequently
\begin{equation}
 \|G\|=\sqrt2,\qquad
 \|G_\uparrow\|=\|G_\downarrow\|=1,\qquad B_0G=d_*G.
 \label{eq:core-spin-components}
\end{equation}
The ionic spatial parity of \(G\) is odd, whereas that of \(\phi\) is
even.  Thus \(QG=G\) at every coupling under consideration.

\begin{lemma}[Cutoff channel equation]
\label{lem:channel-cutoff}
There are an exterior \(\cK\)-valued function \(F\), a nonnegative bounded
multiplication operator \(W(x)\) on the ionic variables, and
\(\sigma>\gamma\) such that
\begin{equation}
 \left(-\frac{\Delta_x}{2}+\frac{k^2}{2}+B+u(x)+\lambda W(x)\right)F=R.
 \label{eq:cutoff-channel-equation}
\end{equation}
The operator \(W\) equals \(\sum_{i=1}^2|x-y_i|^{-1}\) for
\(|y|\leq r/2\), vanishes for \(|y|\geq3r/4\), and satisfies
\begin{equation}
 \|\partial_x^\nu W(x)\|_{\cK\to\cK}
 \leq C_\nu r^{-1-|\nu|}.
 \label{eq:cutoff-interaction-symbol}
\end{equation}
Here \(\nu\) is a Cartesian multiindex.  For every fixed integer \(s\),
\begin{equation}
 \|\partial_x^\nu R(x)\|\leq C_{D,s}\lambda e^{-\sigma r}
 \quad(|\nu|\leq s).
 \label{eq:channel-residual}
\end{equation}
Moreover,
\begin{equation}
 \left|\partial_x^\nu\{n(x)-\|F(x)\|^2\}\right|
 \leq C_D e^{-2\sigma r}\quad(|\nu|\leq2),
 \label{eq:density-cutoff-error}
\end{equation}
and \(F-F_0=O(\lambda)\) on each fixed exterior annulus, also in any
fixed finite number of \(x\)-derivatives.
\end{lemma}

\begin{proof}
Put \(\rho=(|x|^2+|y|^2)^{1/2}\).  A configuration-space partition
according to a particle with distance at least a fixed multiple of
\(\rho\) from the origin gives the exterior form bound
\begin{equation}
 H_{3,\lambda}\geq2\epsilon_b-o(1)
 \quad\hbox{on functions supported where }\rho\to\infty.
 \label{eq:configuration-exterior-bound}
\end{equation}
The nuclear attraction of the selected particle is \(O(1/\rho)\);
the other two one-particle operators are each bounded below by
\(\epsilon_b\), and repulsion can be dropped.  The partition has
derivatives \(O(1/\rho)\), so its kinetic localization error tends
to zero.  This argument is uniform for \(0\leq\lambda\leq1\).

The exponent margin needed below already follows for \(D\ge100\)
from \Cref{eq:two-centre-bottom-lower,eq:two-centre-two-upper}:
\[
 \begin{gathered}
 k_0^2\ge1-\frac{80}{D^2},\qquad
 \gamma^2\le1+\frac2D+\frac{40}{D^2},\\
 \frac{10}{9}k_0^2-\gamma^2
 \ge\frac19-\frac2D-\frac{1160}{9D^2}>0.
 \end{gathered}
\]
We can therefore choose
\[
 0<\tau_0<\tau_1<\tau_2<k_0,\qquad \tau_0\sqrt{10/9}>\gamma.
 \]
After decreasing \(\lambda_D\), the exterior bound and
\(E_3(0)=2\epsilon_b-k_0^2/2\) leave a positive margin above the
weight cost \(\tau_2^2/2\).  The exponential-weight identity, first
with an exterior cutoff and a smooth truncated exponent whose
gradient has norm at most one, gives
\begin{equation}
 \|e^{\tau_1\rho}F_{\mathrm{true}}\|_{H^1(\R^9)}
 +\|e^{\tau_1\rho}F_0\|_{H^1(\R^9)}\leq C_D.
 \label{eq:configuration-weight}
\end{equation}
At zero coupling the same argument also gives
\(\|e^{\tau_2\rho}F_0\|_{H^1}\leq C_D\), since \(\tau_2<k_0\).
The finite spin sum is understood.  The estimate also applies to an
inhomogeneous equation: if
\((H_{3,\lambda}-E_3(\lambda))v=f\), then, for a fixed ball \(K\),
\begin{equation}
 \|e^{\tau_1\rho}v\|_{H^1}
 \leq C_D\left(\|e^{\tau_1\rho}f\|_2+\|v\|_{H^1(K)}\right).
 \label{eq:inhomogeneous-weight}
\end{equation}
Indeed the exterior spectral lower bound controls the weighted
\(L^2\) norm by the weight identity and Cauchy--Schwarz.  Combining
a small multiple of the kinetic form estimate with that lower bound
also controls the gradient; the strict margin between \(\tau_1\) and
\(\tau_2\) absorbs its additional weight cost.  Hardy's inequality
absorbs the infinitesimally form-bounded Coulomb terms.  These bounds
are independent of the exponent truncation, which can therefore
be removed.

The equation for \(v=F_{\mathrm{true}}-F_0\) has source
\[
 \bigl(E_3(\lambda)-E_3(0)-\lambda V_{\mathrm{ee}}\bigr)F_0.
 \]
Hardy's inequality, with the spare exponent \(\tau_2\) on \(F_0\),
bounds it by \(C_D\lambda\) in weighted \(L^2\).
The compact term is \(O(\lambda)\) by \Cref{prop:small-coupling}.
Thus
\begin{equation}
 \|e^{\tau_1\rho}(F_{\mathrm{true}}-F_0)\|_{H^1}\leq C_D\lambda.
 \label{eq:weighted-state-difference}
\end{equation}

Choose smooth functions of \(|y|/r\), with values in \([0,1]\), such
that \(\chi=1\) on \(|y|\leq r/3\), \(\chi=0\) on \(|y|\geq r/2\),
and \(\zeta=1\) on \(|y|\leq r/2\), \(\zeta=0\) on
\(|y|\geq3r/4\).  Define
\[
 F=F_0+\chi(F_{\mathrm{true}}-F_0),\qquad
 W=\zeta\sum_{i=1}^2|x-y_i|^{-1}.
 \]
On the support of \(\zeta\), \(|x-y_i|\geq r/4\), proving
\Cref{eq:cutoff-interaction-symbol}.  Denote the operator on the
left of \Cref{eq:cutoff-channel-equation} by \(\mathcal L_\lambda\).
Since \(W\) is exact on the support of \(\chi\),
\begin{equation}
 R=(1-\chi)\mathcal L_\lambda F_0
       +[\mathcal L_\lambda,\chi](F_{\mathrm{true}}-F_0),
 \label{eq:cutoff-residual-identity}
\end{equation}
where
\[
 \mathcal L_\lambda F_0
 =\left\{\lambda\left(|y_1-y_2|^{-1}+W\right)
          -E_3(\lambda)+E_3(0)\right\}F_0.
 \]
Thus every residual term is \(O(\lambda)\) and is supported where
\(|y|\geq r/3\).  There,
\begin{equation}
 \rho\geq\sqrt{10/9}\,r.
 \label{eq:cutoff-geometric-margin}
\end{equation}

We give the derivative details, since higher ordinary derivatives
of the uncut wavefunction across a collision cannot simply be
iterated.  Local \(H^2\) estimates on nested fixed-size configuration
boxes have uniform constants: the Laplacian estimate with cutoffs
and the infinitesimal Coulomb bound absorb the potential terms;
the local energy estimate controls cutoff commutators.  Where
derivatives of \(\chi\) occur, \(x\) is separated from both other
electrons.  Further \(x\)-derivatives of the equation differentiate
only smooth bounded coefficients, while singular ionic coefficients
are independent of \(x\).  Difference quotients followed by these
local estimates give any fixed number of \(x\)-derivatives of the
difference, at size \(O(\lambda)\) in weighted local norms.
In the other residual term, \(F_0\) has the product form
\Cref{eq:unperturbed-fibers}; its \(x\)-derivatives have no
electron-electron singularity, and the ionic Coulomb multiplier
is controlled by Hardy's inequality.  Summing the box estimates
and using Sobolev embedding in the three \(x\)-coordinates gives
\Cref{eq:channel-residual}.  A small loss from \(\tau_1\) toward \(\tau_0\)
absorbs polynomial integration factors.  By
\Cref{eq:cutoff-geometric-margin}, one still has \(\sigma>\gamma\).

Density derivatives in the discarded region require a different
coordinate choice.  Choose smooth pair cutoffs \(\chi_{ij}\in[0,1]\)
equal to one at distances at most \(1\) and zero at distances at least
\(2\), and expand the product of the three identities
\(1=\chi_{ij}+(1-\chi_{ij})\).
For each term use the connected components of its near-edge graph
as clusters.  They have diameter at most \(4\), and distinct clusters
are separated by at least \(1\).  Nested boxes in which each
particle moves by less than \(1/8\) preserve a positive separation.
The cluster containing \(x\) is far from the nuclei when \(r\) is large.
Translate with \(x\) all other
variables in that cluster.  A derivative in the translated
\(x\)-coordinate is the sum of the corresponding particle
derivatives.  It annihilates within-cluster Coulomb singularities
and singularities involving only outside particles; every
coefficient differentiated is smooth with uniformly bounded
derivatives.  Nested-box energy and elliptic estimates can
therefore be iterated in these three translation directions.
This is the parallel-derivative method of
\cite[preprint version, Proposition~1.10 and Section~6]{FS18}; here the fixed
separations make the constants uniform, while
\Cref{eq:configuration-weight,eq:weighted-state-difference}
preserve the requisite exponential weights.

The difference \(n-\|F\|^2\) is a sum of integrals of products of
\(F_{\mathrm{true}}\) and \(F_0\), with coefficients supported where
\(|y|\geq r/3\).  After the coordinate changes just described,
differentiate these integrals and apply Cauchy--Schwarz to the
weighted local estimates.  Each product has two exponentially
decaying factors.  Sobolev embedding in the three translation
directions and \Cref{eq:cutoff-geometric-margin} give
\Cref{eq:density-cutoff-error}, after another arbitrarily small
loss of exponent.  To obtain smoothness on a compact set disjoint from
the nuclei, replace the pair-cutoff radii \(1,2\) by
\(\varepsilon,2\varepsilon\), with \(\varepsilon\) smaller than one
sixteenth of the set's distance from the nuclei.  The cluster containing
\(x\) has diameter at most \(4\varepsilon\); sufficiently small nested
boxes therefore keep it away from the nuclei and preserve the
intercluster separation.  The same argument without discarding a tail
then proves smoothness of the density there.
On fixed exterior annuli the differentiated local estimates also
give the claimed \(O(\lambda)\) convergence of \(F-F_0\).
\end{proof}

\subsection{Localized multipoles and the slow channel}

Reduced ionic inverses preserve the moment bounds needed below.
Let \(\rho_y=(1+|y|^2)^{1/2}\).  For some \(s_D>0\),
\(\phi\) is bounded uniformly in \(e^{s_D\rho_y}H^2_y\).
Moreover, if \(e^{s\rho_y}v\in L^2\), then
\begin{equation}
 e^{s'\rho_y}B_Q^{-1}Qv\in H^2_y
 \quad\hbox{for some }0<s'<s,
 \label{eq:weighted-ionic-inverse}
\end{equation}
with uniform bounds on the small coupling interval.
To justify this, conjugate the ionic operator by a small exponential
with smooth truncated exponent.  The added first-order term has
small relative operator bound and the added zeroth-order term is
bounded and small.  An isolated contour around the ground
eigenvalue has uniformly bounded conjugated resolvents.  Applying
the resulting projection to a compactly supported vector with
nonzero ground-state overlap proves localization of \(\phi\).
The operator \(B+P\) has a bounded inverse into \(H^2_y\).
Its conjugated finite-rank term is now bounded by the localization
of \(\phi\), so the same inverse perturbation argument applies.
Restriction to \(Q\) proves \Cref{eq:weighted-ionic-inverse}.
Bounds are independent of the exponent truncation.  Polynomial
multiplication is harmless after decreasing the exponent again;
only finitely many such decreases will be used.

On these localized vectors, finite Taylor expansion gives
\begin{equation}
 W=\frac2r+\frac{W_2(\omega)}{r^2}
       +\frac{W_3(\omega)}{r^3}+\cdots,\qquad
 W_2(\omega)=\omega\cdot(y_1+y_2).
 \label{eq:ionic-multipoles}
\end{equation}
An expansion through \(r^{-m}\), applied to a vector bounded in
one of the preceding weighted spaces, has remainder \(O(r^{-m-1})\),
with Cartesian derivative of order \(|\nu|\) bounded by
\(O(r^{-m-1-|\nu|})\).
On \(|y|\leq3r/4\), the Taylor remainder is bounded by a constant
times \(|y|^m r^{-m-1}\).  Outside this region the polynomial
moments of the localized vector bound the omitted cutoff tail
exponentially.  Differentiation gives the same symbol bounds.
These are expansions on localized vectors, not operator-norm
expansions on all of \(\cK\).  Similarly write
\[
 u=-2/r+\sum_{j\geq2}U_j(\omega)r^{-j},\qquad
 U_2=0,\qquad U_3=-D^2(3\omega_z^2-1),
 \]
with finite-order Taylor remainders understood.

\begin{lemma}[Slow-channel expansion]
\label{lem:channel-expansion}
Write \(PF=\phi f\).  With every fixed finite number of angular
derivatives,
\begin{equation}
 \|F\|=O(A),\qquad \|QF\|=O(A/r^2).
 \label{eq:channel-leading-bounds}
\end{equation}
There are smooth real odd scalar functions \(a_0,a_1,a_2\) on
\(\mathbb S^2\) such that
\begin{equation}
 f=A\left(a_0+\frac{a_1}{r}+\frac{a_2}{r^2}+\mathcal E\right).
 \label{eq:projected-channel-expansion}
\end{equation}
For each fixed angular differential operator \(\Omega^\ell\),
\begin{equation}
 \Omega^\ell\mathcal E=O(r^{-3}),\qquad
 \partial_r\Omega^\ell\mathcal E=O(r^{-4}),\qquad
 \partial_r^2\Omega^\ell\mathcal E=O(r^{-4}).
 \label{eq:scalar-remainder-derivatives}
\end{equation}
The radial remainders are odd in \(\omega_z\); in particular their
bounds without angular differentiation gain a factor
\(|\omega_z|\), also after one or two radial derivatives.
All bounds are uniform for small \(\lambda\).
Furthermore \(a_0^\lambda\to a_0^0\) in \(C^1(\mathbb S^2)\), and
phases can be chosen so that
\begin{equation}
 \inf_{\omega_z>0}\frac{a_0^\lambda(\omega)}{\omega_z}>c_D>0.
 \label{eq:simple-angular-zero}
\end{equation}
Finally,
\begin{equation}
 b=e^{-\gamma r}r^{2/\gamma-1}
       \bigl(b_\infty(\omega)+O(r^{-1})\bigr),\qquad
 \min_{\mathbb S^2}b_\infty>0,
 \label{eq:fast-orbital-expansion}
\end{equation}
with angular derivatives and two radial derivatives of the expansion.
\end{lemma}

\begin{proof}
The ionic gap and multipole estimate imply
\[
 PW\phi=(2/r+O(r^{-2}))\phi,\qquad QW\phi=O(r^{-2}).
 \]
Splitting a vector into \(P+Q\) and applying Young's inequality to
the off-diagonal term, while retaining half the positive \(Q\)-gap,
gives
\begin{equation}
 B+\lambda W\geq 2\lambda/r-C_D/r^2.
 \label{eq:ionic-schur-bound}
\end{equation}
The norm \(N(x)=\|F(x)\|\) consequently satisfies
\[
 \left(-\Delta_x/2+k^2/2-(2-2\lambda)/r-C_D/r^2\right)N\leq\|R\|.
 \]
Take the real inner product of the vector equation with \(F/\|F\|\)
and use \(|\nabla\|F\||\leq\|\nabla F\|\).  Replacing the norm by
\((\|F\|^2+\varepsilon^2)^{1/2}\) proves the weak inequality also
at zeros.  The unbounded ionic operator is handled through its
quadratic form, or by bounded spectral truncations.

Applying this scalar operator to \(A(1-M/r)\), and dividing by \(A\),
gives \(r^{-2}\) coefficient
\[
 kM-\beta(\beta+1)/2-C_D.
 \]
Choose \(M\) to make this positive, and start outside \(r>2M\).
A sufficiently large multiple dominates the inner boundary value
and the exponentially smaller source, yielding \(N=O(A)\) by
exterior comparison.  The latter can be proved by testing the
positive part of the difference with cutoffs: the zeroth-order
coefficient is positive sufficiently far out, both functions
decay, and the cutoff error tends to zero.
The \(Q\)-equation has source
\(-\lambda QW\phi f+QR=O(A/r^2)\).  Its norm inequality retains
the positive constant \(c_D\) in addition to \(k^2/2+u\), so
\(A/r^2\) is a supersolution at sufficiently large radius.  This
proves \Cref{eq:channel-leading-bounds}.

Rotation fields \(\Omega\) acting on \(x\) commute with
\(\Delta_x,B,P,Q\).  Although \(\|\Omega^jW\|=O(r^{-1})\),
one has \(\|\Omega^jW\phi\|=O(r^{-2})\) for \(j\geq1\), and
\(\Omega^ju=O_D(r^{-3})\).
Suppose the two leading bounds hold for fewer than \(s\) rotations.
Every commutator source for \(\Omega^sF\) is \(O(A/r^2)\):
its \(P\)-part uses the improved estimate on \(W\phi\), and its
\(Q\)-part uses \(O(r^{-1})O(A/r^2)\).
First the full norm comparison gives \(\Omega^sF=O(A)\);
then the \(Q\)-equation and its gap give \(Q\Omega^sF=O(A/r^2)\).
This closes the induction.  The absolute decay needed for
comparison follows from the cutoff construction, allowing a small
exponential loss for the polynomial coefficients of rotation fields.

Set \(h_0=f/A\), and define
\begin{equation}
 K_j=-\tfrac12\Delta_{\mathbb S^2}
       -\tfrac12(\beta-j)(\beta-j+1),\qquad
 c_m=U_m+\lambda\langle\phi,W_m\phi\rangle.
 \label{eq:angular-recursion-operators}
\end{equation}
Axial symmetry and reflection give \(c_2=0\).
The projected equation conjugated by \(A\) is
\begin{equation}
 -\tfrac12h_0''+
 \left(k-\frac{\beta+1}{r}\right)h_0'
 +\frac{K_0+c_2}{r^2}h_0+\frac{c_3}{r^3}h_0+\cdots
 =O(r^{-4}),
 \label{eq:projected-radial-equation}
\end{equation}
with any fixed number of angular derivatives.
The off-diagonal source has this order because
\(\|PWQ\|=\|QW\phi\|=O(r^{-2})\) and \(QF=O(A/r^2)\).

For the radial step, if a bounded angular-valued function \(H\)
satisfies
\[
 -\tfrac12H''+\left(k-\frac{\beta+1}{r}\right)H'=G(r),
 \qquad G(r)=O(r^{-q}),\quad q>1,
 \]
then
\begin{equation}
 H'(r)=2e^{2kr}r^{-2(\beta+1)}
       \int_r^\infty e^{-2ks}s^{2(\beta+1)}G(s)\,ds.
 \label{eq:radial-integrating-factor}
\end{equation}
The other homogeneous branch grows exponentially and is excluded
by boundedness.  The integral is \(O(r^{-q})\), uniformly for the
present bounded \(\beta\) and positive lower bound on \(k\);
thus \(H\) has a limit with error \(O(r^{1-q})\).
Initially the right side of the radial equation is \(O(r^{-2})\),
since all needed angular derivatives of \(h_0\) are bounded.
Hence \(h_0\to a_0\), with error \(O(r^{-1})\).
Set
\begin{equation}
 a_1=\frac{(K_0+c_2)a_0}{k},\qquad
 a_2=\frac{(K_1+c_2)a_1+c_3a_0}{2k}.
 \label{eq:scalar-channel-recursion}
\end{equation}
Subtracting \(a_0+a_1/r\) cancels order two; the integrating factor
improves the remainder to \(O(r^{-2})\), with derivative
\(O(r^{-3})\).  Subtracting \(a_2/r^2\) cancels order three,
giving \(\mathcal E=O(r^{-3})\) and \(\mathcal E'=O(r^{-4})\).
The equation gives \(\mathcal E''=O(r^{-4})\).
The same argument with finitely more angular derivatives at each
step proves \Cref{eq:scalar-remainder-derivatives}.
Reflection makes \(f,a_j\), and all radial remainders odd.
The mean-value theorem in the transverse angular coordinate gives
the additional factor at the equator.

On a fixed radius, the projected function and its angular derivatives
converge by the local estimates in \Cref{lem:channel-cutoff}.
Since \(\|f/A-a_0\|_{C^1}\leq C_D/r\) uniformly, first choosing
the radius large and then sending \(\lambda\) to zero proves
\(C^1\) convergence of \(a_0^\lambda\).
At zero coupling the scalar equation for \(e\) on the upper half-space
has a subsolution which is a positive multiple of
\[
 e^{-k_0r}r^{2/k_0-1}\omega_z(1+M/r).
 \]
After division by the radial factor and \(\omega_z\), the controlling
\(r^{-2}\) coefficient contains \(-k_0M\); it dominates the bounded
angular and scalar corrections for large \(M\).
Positivity of \(e\) on an inner hemisphere and the boundary point
lemma on its equator supply a strictly positive starting multiple,
including the equatorial limit.  Exterior comparison proves
\(\inf a_0^0/\omega_z>0\).  Oddness and \(C^1\) convergence give
\Cref{eq:simple-angular-zero}.  The identical scalar argument for
\(b\), without an angular zero and with rate \(\gamma\), proves
\Cref{eq:fast-orbital-expansion}; positivity on the compact inner
sphere gives \(\min b_\infty>0\).
\end{proof}

\subsection{The first nonzero complementary term on the plane}

The scalar channel's oddness forces it to vanish on the median plane.
We now identify the complementary term that supplies its missing
density there, and control its error before comparing it with the
retained fast term \(bG\).

\begin{lemma}[Complementary expansion and plane coefficient]
\label{lem:plane-coefficient}
There are smooth exponentially localized \(Q\cK\)-valued angular
coefficients \(T_2,T_3,T_4\), bounded uniformly for small \(\lambda\),
such that
\begin{equation}
 QF=bG+\lambda A
       \left(\frac{T_2}{r^2}+\frac{T_3}{r^3}+\frac{T_4}{r^4}\right)
       +\mathcal E_Q,\qquad
 \|\partial_x^\nu\mathcal E_Q\|
       \leq C_D\lambda Ar^{-5}\quad(|\nu|\leq2).
 \label{eq:complementary-expansion}
\end{equation}
Writing \(Y_z=y_{1z}+y_{2z}\), one has
\begin{equation}
 \begin{split}
 T_2&=-a_0B_Q^{-1}QW_2\phi,\\
 T_2|_{\omega_z=0}&=T_3|_{\omega_z=0}=0,\\
 T_4|_{\omega_z=0}
   &=\frac12B_Q^{-1}\Delta_{\mathbb S^2}T_2|_{\omega_z=0}
     =-a_0'(0)B_Q^{-2}QY_z\phi .
 \end{split}
 \label{eq:equatorial-coefficient}
\end{equation}
The last vector has norm bounded below by a positive constant at
fixed \(D\), uniformly for sufficiently small \(\lambda\).
All \(T_j\) have distinguished spin up.
\end{lemma}

\begin{proof}
On localized angular coefficients define
\[
 C_m=U_m\,\mathrm{id}_{Q\cK}+\lambda QW_mQ.
 \]
The qualification is essential: \(W_m\), and hence \(C_m\), can
be unbounded on arbitrary ionic vectors.  Their action on the
localized vectors in this proof is controlled by
\Cref{eq:weighted-ionic-inverse,eq:ionic-multipoles}.
On a term \(r^{-j}T(\omega)\), the conjugated \(Q\)-operator has
expansion
\begin{equation}
 r^{-j}B_QT-jk r^{-j-1}T
 +r^{-j-2}K_jT+\sum_{m\geq2}r^{-j-m}C_mT.
 \label{eq:complementary-recursion-action}
\end{equation}
The source from \(P\), after division by \(\lambda A\), is
\(-QW\phi\,f/A\).  Its coefficients of orders two through four are
\[
 \begin{split}
 S_2&=-a_0QW_2\phi,\\
 S_3&=-a_1QW_2\phi-a_0QW_3\phi,\\
 S_4&=-a_2QW_2\phi-a_1QW_3\phi-a_0QW_4\phi.
 \end{split}
 \]
We therefore define the coefficients by
\begin{equation}
 \begin{split}
 B_QT_2&=S_2,\\
 B_QT_3&=S_3+2kT_2,\\
 B_QT_4&=S_4+3kT_3-(K_2+C_2)T_2.
 \end{split}
 \label{eq:complementary-recursion}
\end{equation}
Repeated use of \Cref{eq:weighted-ionic-inverse}, decreasing a
positive weight after polynomial multiplication when necessary,
gives smooth localized coefficients with every finite angular
derivative bound used here.  The nuclear quadrupole has already
entered \(a_2\) in \Cref{eq:scalar-channel-recursion};
the operator \(C_3\) acts directly on this complementary expansion
first at order five.

Every \(S_j\) vanishes on the equator because its scalar factors
\(a_i\) are odd.  The first two recursion equations give
\(T_2=T_3=0\) there.  In the third equation all terms except the
angular Laplacian of \(T_2\) also vanish.  At the equator
\(a_0=\Delta_{\mathbb S^2}a_0=0\), while
\[
 \nabla_{\mathbb S^2}a_0=a_0'(0)e_z,\qquad
 \partial_{\omega_z}W_2=Y_z.
 \]
The product rule in \(T_2=-a_0B_Q^{-1}QW_2\phi\) proves
\Cref{eq:equatorial-coefficient}, with its stated sign and factor.
The slope is nonzero by \Cref{eq:simple-angular-zero}.
Moreover \(Y_z\phi\) is ionic-odd, so \(QY_z\phi=Y_z\phi\).
This vector cannot vanish: the zero set of the multiplier \(Y_z\)
is a measure-zero hyperplane, which cannot support a nonzero
\(L^2\) vector.  Injectivity of \(B_Q^{-2}\) proves nonvanishing.
Continuity of the localized resolvents and vectors gives the
uniform lower bound for small \(\lambda\).  All operators in the
recursion preserve the distinguished spin, so each \(T_j\) has
spin up.

We next prove the remainder with its asserted uniformity.
The retained vector \(bG\) solves the homogeneous \(Q\)-equation
at \(\lambda=0\).  Since \(QG=G\), its perturbed residual is
\begin{equation}
 \begin{split}
 \mathcal L_{\lambda,Q}(bG)
 &=bQ\left\{\lambda\left(|y_1-y_2|^{-1}+W\right)
                 -E_3(\lambda)+E_3(0)\right\}G\\
 &=O(\lambda b),
 \end{split}
 \label{eq:fast-channel-residual}
\end{equation}
also through two ordinary \(x\)-derivatives.
Hardy's inequality controls the ionic Coulomb factor on \(G\);
the other \(x\)-dependent factors are smooth in the exterior.
By \Cref{eq:channel-rate-separation,eq:fast-orbital-expansion},
\begin{equation}
 b/A\leq C_D e^{-(\gamma-k_0)r/2}r^{C_D}.
 \label{eq:fast-term-absorption}
\end{equation}
Thus the residual is bounded by \(C_D\lambda Ar^{-5}\), uniformly
in \(\lambda\).  No equality of perturbed and unperturbed excited
thresholds is being assumed.

The matched slow terms have residual \(O(\lambda Ar^{-5})\).
The same bound holds through two Cartesian derivatives:
finite multipole remainders have symbol bounds on the localized
coefficients, and the scalar remainder is multiplied by
\(\lambda QW\phi=O(\lambda r^{-2})\).
The three bounds in \Cref{eq:scalar-remainder-derivatives} suffice;
one does not need the stronger assertion
\(\partial_r^2\mathcal E=O(r^{-5})\).
The source \(QR\) is exponentially smaller.
Hence \(\mathcal E_Q\) satisfies a \(Q\)-equation with a source
of this size, including two derivatives.  On a fixed large
starting sphere its value is \(O(\lambda)\), since \(QF=bG\)
at zero coupling and all fixed-radius data converge at that rate.
The positive \(Q\)-gap makes a sufficiently large multiple of
\(\lambda Ar^{-5}\) a scalar norm supersolution which dominates
both the source and boundary value.  Exterior comparison gives
the undifferentiated remainder bound.

For the derivative upgrade it is important to include the unbounded
ionic operator.  For compactly supported \(v\) with values in
\(\operatorname{Dom}B_Q\), Fourier transform in \(x\) and the
spectral resolution of \(B_Q\geq0\) give
\begin{equation}
 \|D_x^2v\|_2+\|B_Qv\|_2
 \leq C\|(-\Delta_x/2+B_Q)v\|_2.
 \label{eq:positive-ionic-elliptic-estimate}
\end{equation}
Indeed \(|\xi|^2\) and a spectral value \(b\geq0\) are each bounded
by a fixed multiple of \(|\xi|^2/2+b\).
For the local estimate, first control \(\nabla_xv\) by the
cutoff energy identity and then apply
\Cref{eq:positive-ionic-elliptic-estimate} with a second cutoff.
The remaining coefficient \(k^2/2+u+\lambda QWQ\) and its
\(x\)-derivatives are bounded operators on exterior unit balls,
uniformly in their centers.  Since \(B_Q\) is independent of \(x\),
difference quotients and two further local estimates give
\[
 \|\mathcal E_Q\|_{H_x^4(B(x,1/2);\cK)}
 \leq C_D\lambda A(r)r^{-5}.
 \]
The source derivatives and the already proved pointwise bound
control the right side on a larger ball; \(A\) varies by a bounded
factor there.  Hilbert-valued Sobolev embedding in dimension three,
\(H^4\hookrightarrow C^{2,\alpha}\) for \(\alpha<1/2\), gives
the two pointwise derivative bounds in
\Cref{eq:complementary-expansion}.
\end{proof}

\subsection{Uniform comparison of the three density scales}

The coefficient calculation concerns the eventual tail at a fixed
positive coupling.  To obtain a common exterior radius as the coupling
tends to zero, we must also control the region where the fast term and
the induced plane term are comparable.  Keeping all three scales in
the following lower bound makes that comparison uniform.

\begin{proof}[Proof of \Cref{prop:density-properties}]
Put \(t=\omega_z\) and
\[
 T_{\mathrm{eq}}=T_4|_{t=0},\qquad
 M(x)=\phi f(x)+b(x)G+\lambda A(r)r^{-4}T_{\mathrm{eq}}.
 \]
Axial symmetry and \Cref{eq:equatorial-coefficient} make
\(T_{\mathrm{eq}}\) independent of equatorial azimuth.
Smoothness and the zeros of \(T_2,T_3\) give
\[
 \|T_2(\omega)\|+\|T_3(\omega)\|
       +\|T_4(\omega)-T_{\mathrm{eq}}\|\leq C_D|t|.
 \]
The estimate holds on the whole sphere by compactness away from
the equator.  For any distinguished-spin-up \(Q\)-vector \(S\),
\Cref{eq:core-spin-components} implies
\[
 |\langle bG,S\rangle|
 \leq |b|\,\|S\|=2^{-1/2}\|bG\|\,\|S\|.
 \]
Therefore
\begin{equation}
 \begin{split}
 (1-2^{-1/2})(2b^2+\|S\|^2)
 &\leq\|bG+S\|^2\\
 &\leq(1+2^{-1/2})(2b^2+\|S\|^2).
 \end{split}
 \label{eq:spin-noncancellation}
\end{equation}
The \(P\)-component is orthogonal to both these terms.  The simple
angular zero and positive lower bound on \(\|T_{\mathrm{eq}}\|\)
give
\[
 \|M\|\geq c_D(A|t|+b+\lambda Ar^{-4}).
 \]
The complementary expansion gives
\[
 \|F-M\|\leq C_D\lambda A(r^{-2}|t|+r^{-5})=o_D(\|M\|),
 \]
uniformly in small \(\lambda\) as the radius increases.
Consequently
\begin{equation}
 \|F\|\geq c_D(A|\omega_z|+b+\lambda Ar^{-4}).
 \label{eq:three-scale-lower-bound}
\end{equation}
In addition, there is the absolute bound
\begin{equation}
 b/\|F\|\leq2.
 \label{eq:absolute-core-ratio}
\end{equation}
The distinguished-down-spin component of \(M\) has norm \(b\).
Increase the radius so that \(\|F-M\|\leq\|M\|/2\); then
\[
 \|F\|\geq\|M\|-\|F-M\|\geq\|M\|/2\geq b/2.
\]
Only the radius, not the constant \(2\),
depends on \(D\).
The cutoff errors are relatively negligible because
\[
 e^{-2\sigma r}/b^2\longrightarrow0\qquad(\sigma>\gamma).
 \]
This proves exterior positivity of \(n\); positivity on the
remaining compact set follows from
\Cref{lem:density-regularity} after decreasing \(\lambda_D\).
Symmetry, compact convergence, and exterior smoothness have
already been proved.

We verify differentiated comparisons through the moving crossover.
The projected remainders and their first two radial derivatives
are odd, so their angular derivative bounds supply the factor
\(|t|\).  Their errors can thus be divided by the first lower
scale in \Cref{eq:three-scale-lower-bound}.
For the complementary error,
\[
 \frac{\|\partial_x^\nu\mathcal E_Q\|}{\|F\|}
 \leq C_D/r\quad(|\nu|\leq2,\ \lambda>0),
 \]
using the third lower scale.  At \(\lambda=0\) this error is zero.
Differentiating the explicit slow terms with \(\partial_r+k\)
leaves only inverse powers of \(r\).
For the fast orbital,
\[
 (\partial_r+k)b=(k-\gamma)b+O_D(b/r),\qquad
 (\partial_r+k)^2b=(k-\gamma)^2b+O_D(b/r).
 \]
It follows, with \(\delta_\lambda=k-\gamma\), that
\begin{equation}
 \begin{split}
 (\partial_r+k)F&=\delta_\lambda bG+o_D(\|F\|),\\
 (\partial_r+k)^2F&=\delta_\lambda^2bG+o_D(\|F\|),
 \end{split}
 \label{eq:radial-relative-derivatives}
\end{equation}
where the ratios denoted by \(o_D(1)\) tend to zero uniformly
in the chosen coupling interval.

At fixed \(z\), \(\partial_p\) has radial part
\((p/r)\partial_r\) and angular part of size \(O(|t|/r)\)
on \(0\leq z\leq p/2\).
Applying it once or twice to smooth odd angular coefficients
retains a factor \(|t|\), together with the corresponding
inverse powers of \(r\).  Thus
\[
 \left(\partial_p+k\frac p r\right)F
 =\delta_\lambda\frac p r\,bG+o_D(\|F\|).
 \]
On \(0\leq z\leq L\), \(1-p/r=O_L(p^{-2})\) is negligible.
With \(D_p=\partial_p+k\),
\begin{equation}
 D_pF=\delta_\lambda bG+o_D(\|F\|),\qquad
 D_p^2F=\delta_\lambda^2bG+o_D(\|F\|).
 \label{eq:horizontal-relative-derivatives}
\end{equation}
These estimates include angular derivatives of the displayed
\(T_j\): their zeros are retained for \(T_2,T_3\), whereas the
extra angular factor for \(T_4\) is bounded using the first lower
scale.  The ordinary derivative bounds handle \(\mathcal E_Q\)
directly.

By \Cref{eq:absolute-core-ratio}, \(\|bG\|/\|F\|\leq2\sqrt2\).
The identity
\[
 (\partial_r+2k)\|F\|^2
 =2\operatorname{Re}\langle F,(\partial_r+k)F\rangle
 \]
and its horizontal version bound the normalized first derivatives
by \(4\sqrt2|\delta_\lambda|+o_D(1)\).
Also
\[
 (\partial_p+2k)^2\|F\|^2
 =2\|D_pF\|^2+
  2\operatorname{Re}\langle F,D_p^2F\rangle
 \]
gives the normalized second-derivative bound
\((16+4\sqrt2)\delta_\lambda^2+o_D(1)\) in the strip.
The differentiated cutoff errors remain relatively negligible.
In particular, there is \(\rho_D(R)\downarrow0\), independent
of small \(\lambda\), such that
\begin{equation}
 \begin{split}
 \max\Bigg\{&
 \sup_{r\geq R}|\partial_r\log n+2k|,\ 
 \sup_{\substack{r\geq R\\0\leq z\leq p/2}}
       \left|\partial_p\log n+2k\frac p r\right|,\\
 &\sup_{\substack{p\geq R\\0\leq z\leq L}}
       \frac{|m_p|+|m_{pp}|}{m}\Bigg\}
 \leq100(|\gamma-k|+|\gamma-k|^2)+\rho_D(R).
 \end{split}
 \label{eq:quantitative-density-drift}
\end{equation}
All \(D\)-dependent coefficients other than the fast exponential
rate have entered the vanishing function \(\rho_D\).
Since \(\gamma-k_0\to0\) as \(D\to\infty\) and \(k\to k_0\)
as \(\lambda\downarrow0\), this proves the derivative assertions
in \Cref{eq:density-drift,eq:strip-data}.

For the vertical ratio the expansions give, uniformly on
\(0\leq z\leq L\),
\begin{equation}
 \begin{split}
 F(p,z)={}&\phi a_0'(0)A(p)z/p+b(p,0)G
               +\lambda A(p)p^{-4}T_{\mathrm{eq}}\\
 &+o_D\bigl(A(p)z/p+b(p,0)+\lambda A(p)p^{-4}\bigr).
 \end{split}
 \label{eq:strip-comparison-vector}
\end{equation}
The last coefficient tends to zero uniformly in \(\lambda,z\),
and the error is relatively small compared with the norm of
the three-term comparison vector by
\Cref{eq:spin-noncancellation}.
Indeed \(r-p=O_L(p^{-1})\); after including the lower radial terms,
the projected factor is \(a_0'(0)z/p+O_D(z/p^2)\);
and \(b(p,z)/b(p,0)=1+O_D(p^{-1})\).
The \(T_2,T_3\) terms are \(O_D(\lambda A z/p^3)\);
all other errors have the displayed bound by
\Cref{eq:three-scale-lower-bound}.
The squared comparison norm lies between \(1-1/\sqrt2\) and
\(1+1/\sqrt2\) times
\[
 |a_0'(0)|^2A(p)^2z^2/p^2+
 2b(p,0)^2+\lambda^2A(p)^2p^{-8}\|T_{\mathrm{eq}}\|^2.
 \]
This expression is nondecreasing in \(z\).
The ratio of the constants is \(3+2\sqrt2<6\);
after absorbing uniformly small errors, the ratio in
\Cref{eq:strip-data} is at most \(20\).

Now fix \(\lambda>0\).  On the plane \(f=0\) and \(T_2=T_3=0\).
Since \(b=o(Ap^{-4})\), the complementary expansion yields
\[
 F(p,0)=\lambda A(p)p^{-4}T_{\mathrm{eq}}+o(A(p)p^{-4}).
 \]
The coefficient is nonzero and the density cutoff error is
exponentially smaller, proving \Cref{eq:plane-tail}.

For plane-centered \(x_j=p_j\omega_j\), put \(r_j(x)=|x_j+x|\).
Uniformly on compact displacement sets,
\[
 r_j(x)-p_j\to\omega_*\cdot x,\qquad
 p_jz/r_j(x)\to z,\qquad
 A(r_j(x))/A(p_j)\to e^{-k\omega_*\cdot x}.
 \]
The projected expansion and its simple angular zero imply
\[
 \frac{PF(x_j+x)}{A(p_j)/p_j}
 \longrightarrow\phi a_0'(0)z e^{-k\omega_*\cdot x}.
 \]
The complementary part divided by this normalization tends
uniformly to zero.  Its fast term is exponentially smaller;
the \(T_2,T_3\) terms retain the angular zero and the plane
term is \(O(Ap_j^{-4})\).
The cutoff error is negligible.  Squaring gives
\Cref{eq:plane-translate} with \(a=|a_0'(0)|>0\).

Finally suppose \(r_j=|x_j|\to\infty\) and \(|z_j|\to\infty\).
Reflection reduces the proof to \(z_j>0\).
If \(\omega_{j,z}\) stays bounded below, the conclusion follows
directly from the smooth nonzero projected angular coefficient.
Otherwise \(t_j=z_j/r_j\to0\), and
\[
 \frac{b(r_j)}{A(r_j)t_j}\to0,\qquad
 \frac{\lambda A(r_j)r_j^{-4}}{A(r_j)t_j}
       =\frac{\lambda}{r_j^3z_j}\to0.
 \]
The \(T_2,T_3\) terms are relatively negligible as well.
For bounded \(x\), \(t_j(x)/t_j=1+O(1/z_j+1/r_j)\).
Since \(a_0(t)/t\) is continuous and bounded away from zero
through the equator, the angular amplitude has relative
variation tending to one.  The radial ratio tends to
\(e^{-k\omega_*\cdot x}\); all lower terms are uniformly
negligible.  Squaring proves \Cref{eq:bulk-translate}.
\end{proof}

We record one more zero-coupling consequence for the next section:
\begin{equation}
 \sup_{0<z\leq10}
 \left\{
 \left|\partial_p\log\frac be+(\gamma-k_0)\frac p r\right|
 +\left|\partial_p^2\log\frac be\right|
 \right\}\longrightarrow0\qquad(p\to\infty).
 \label{eq:unperturbed-strip-ratio}
\end{equation}
Indeed the scalar expansions have the strip form
\[
 e(p,z)=z e^{-k_0r}r^{2/k_0-2}c(r,z/r),\qquad
 b(p,z)=e^{-\gamma r}r^{2/\gamma-1}d(r,z/r),
 \]
where \(c,d\) are smooth, bounded above and below by positive
constants, and have the already proved radial and angular
derivative estimates.  The term \(-\log z\) in \(\log(b/e)\)
disappears under \(\partial_p\).  The other terms give
\(\partial_p\log(b/e)=-(\gamma-k_0)p/r+O_D(r^{-1})\) and
\(\partial_p^2\log(b/e)=O_D(r^{-2})+O_D(z^2/r^3)\).
This proves the assertion uniformly down to the simple zero of \(e\).

\subsection{A fixed-parameter version}

The exact charge will be defined by forward spectral and variational
tests.  Passing them need not put its reciprocal in a particular
perturbative interval chosen earlier.  The next proposition isolates
the hypotheses needed for the final tail argument.

\begin{proposition}[Channels at a fixed coupling]
\label{prop:fixed-parameter-channels}
Suppose \(D\) satisfies \Cref{prop:two-centre-spectrum} and
\[
 k_0>\sqrt{0.98},\qquad \gamma<\sqrt{49/45}.
 \]
Fix \(t\in(0,1]\).  Assume the two-electron ground level is isolated
and simple with an even ionic singlet, and the three-electron
ground level in \(S_z=1/2\) is isolated and simple, is an absolute
ground level, and has a real axisymmetric odd ground state.
Assume \(n_t>0\), and set
\[
 k=\sqrt{2(E_2(t)-E_3(t))},\qquad
 A_t(r)=e^{-kr}r^{(2-2t)/k-1}.
 \]
Suppose
\begin{equation}
 \begin{gathered}
 2\epsilon_b-E_3(t)>0.49,\qquad 0.99<k<\gamma,\\
 \inf_{\omega_z>0}\liminf_{r\to\infty}
 \frac{\|P_tF_{\mathrm{true},t}(r\omega)\|}
      {\omega_z A_t(r)}>0.
 \end{gathered}
 \label{eq:fixed-channel-hypotheses}
\end{equation}
Then \(n_t\) has the local regularity in
\Cref{lem:density-regularity}, is smooth away from the nuclei,
and satisfies \Cref{eq:plane-tail,eq:plane-translate,eq:bulk-translate}
at this fixed coupling, with \(A=A_t\) and suitable positive constants.
Membership of \(t\) in the preceding perturbative interval is
not required.
\end{proposition}

\begin{proof}
Local regularity uses only the Coulomb eigenvalue equation, and
exterior smoothness follows from the cluster argument in
\Cref{lem:channel-cutoff}.  The isolated ionic ground supplies
a positive \(Q\)-gap and \Cref{eq:weighted-ionic-inverse},
with constants now allowed to depend on \(t\).

The configuration-space lower bound
\Cref{eq:configuration-exterior-bound} remains valid because
\(t\geq0\).  The energy hypothesis permits weights with exponent
arbitrarily close to \(\sqrt{0.98}\).  Since
\[
 \sqrt{0.98}\sqrt{10/9}=\sqrt{49/45}>\gamma
 \]
and \(k_0>\sqrt{0.98}\), the true and unperturbed states admit
the same cutoff construction with residual exponent
\(\sigma>\gamma\).
There is no need to estimate the difference uniformly as
\(t\downarrow0\): at this fixed positive \(t\), its finite
weighted norms and the fixed energy difference divided by \(t\)
are part of the constants.  Thus
\Cref{eq:cutoff-channel-equation,eq:channel-residual,eq:density-cutoff-error}
hold at the chosen parameter.

The norm comparisons, angular induction, and integrating-factor
argument use only the positive ionic gap, bounded exterior
coefficients, and exponential residual.  They give
\Cref{eq:projected-channel-expansion,eq:scalar-remainder-derivatives}
without a perturbative assertion about \(a_0\).
Also \(\|F_{\mathrm{true}}-F\|=O(e^{-\sigma r})\), by the same
weighted local estimates on the tail support.
For each fixed \(\omega_z>0\), the liminf in
\Cref{eq:fixed-channel-hypotheses} is therefore
\(|a_0(\omega)|/\omega_z\).
Its positive infimum and the smooth odd extension imply
\[
 \inf_{\omega_z>0}\frac{|a_0(\omega)|}{\omega_z}>0,\qquad
 |a_0'(0)|>0.
 \]
The sign of \(a_0\) is constant on the upper hemisphere and does
not affect the density.

Repeat the complementary recursion
\Cref{eq:complementary-recursion}.  The strict inequality
\(k<\gamma\) absorbs the residual of \(bG\) into \(tAr^{-5}\);
the stronger uniform separation in
\Cref{eq:channel-rate-separation} is unnecessary.
On a fixed starting sphere the remainder is finite; dividing its
bound by the positive profile \(tA_t(r)r^{-5}\) is absorbed in the
allowed \(t\)-dependent comparison constant.
The norm and differentiated comparisons give
\Cref{eq:complementary-expansion} with \(t\)-dependent constants.
Finally \(QY_z\phi=Y_z\phi\ne0\), so the coefficient in
\Cref{eq:equatorial-coefficient} is nonzero without continuity
from zero coupling.  The fixed-coupling parts of the preceding
density proof give the plane lower bound and both translation limits.
\end{proof}
